% !TeX root = ../main.tex
\chapter{Reproducibility and Supplementary Materials}
\label{app:reproducibility}
\thesisplaceholder{Document reproducibility materials and supporting analyses using actual records.}

\section{Data Dictionary and Cohort Definitions}
% 列明公開允許的欄位名稱、單位、取值、缺失定義、時間窗及納入排除規則。
% 不放病人識別資料、受限原始記錄或可還原個體的交叉表。

\section{Implementation and Environment}
% 記錄程式版本、環境、套件、設定檔、隨機種子與訓練／評估指令。
% 若資料不可公開，說明合法取得流程及可公開的程式或合成測試資料。

\section{Additional Analyses}
% 補充分析須對應正文問題，標明預先規劃或探索性。
% 說明各圖表來源、版本與統計單位，勿把探索性結果當確認性證據。

\section{Reporting and Availability}
% 說明採用的報告指引及適用理由；僅標示實際完成的項目。
% 列出可提供與不可提供的材料、限制原因及保存位置；不承諾未獲授權的分享。
